\subsection{Differential properties of finitely generated extensions}

THEOREM 4. --- Let P be a perfect field, L an extension of P and K a subextension of L; we suppose that L is a finitely generated extension of K. Then the canonical L-linear mapping \(\alpha: \Omega_{\mathrm{P}}(\mathrm{K}) \otimes_{\mathrm{K}} \mathrm{L} \to \Omega_{\mathrm{P}}(\mathrm{L})\) has index equal to the transcendence degree of L over K.

If E and F are two subextensions of L such that \(E \subset F\) , we denote by \(\alpha(F/E)\) the canonical F-linear mapping of \(\Omega_{\mathrm{P}}(\mathrm{E}) \otimes_{\mathrm{E}} F\) into \(\Omega_{\mathrm{P}}(\mathrm{F})\) and when it is defined, the index of \(\alpha(F/E)\) will be written \(d(F/E)\) . If E, F and G are three subextensions of L such that \(E \subset F \subset G\) , we have a commutative diagram

\[
\begin{array}{c} \Omega_ {\mathrm{P}} (\mathrm{F}) \otimes_ {\mathrm{F}} \mathrm{G} \xrightarrow {\alpha (\mathrm{G} / \mathrm{F})} \Omega_ {\mathrm{P}} (\mathrm{G}) \\ \alpha (\mathrm{F} / \mathrm{E}) \otimes_ {\mathrm{F}} \mathrm{Id} _ {\mathrm{G}} \Bigg | _ {\uparrow} \\ (\Omega_ {\mathrm{P}} (\mathrm{E}) \otimes_ {\mathrm{E}} \mathrm{F}) \otimes_ {\mathrm{F}} \mathrm{G} \xrightarrow {\beta} \Omega_ {\mathrm{P}} (\mathrm{E}) \otimes_ {\mathrm{E}} \mathrm{G} \end{array}
\]

where \(\beta\) is the canonical isomorphism defined in Prop. 2 (II, p.~278). Since the index is clearly invariant under extension by scalars and the index of an isomorphism is zero, Lemma 2 (V, p.~133) shows that the index \(d(\mathrm{G}/\mathrm{E})\) is defined when \(d(\mathrm{F}/\mathrm{E})\) and \(d(\mathrm{G}/\mathrm{F})\) are and then

\[
d (\mathrm{G} / \mathrm{E}) = d (\mathrm{G} / \mathrm{F}) + d (\mathrm{F} / \mathrm{E}).\tag{5}
\]

Since L is a finitely generated extension of K, Formula (5) and the Cor. of V, p.~111 shows that it is enough to consider the case where x exists such that \(L = K(x)\) ; further if x is algebraic over K, there exists a power q of the characteristic exponent of L such that \(x^{q}\) is separable algebraic over K (V, p.~44, Prop. 13). So it is enough to establish the equality \(d(\mathrm{L}/\mathrm{K}) = \operatorname{tr} \cdot \deg_{\mathrm{K}} \mathrm{L}\) in the three special cases below:

\begin{enumerate}
\def\labelenumi{\alph{enumi})}
\item
  \(x\) is transcendental over \(\mathbf{K}\) : then \(\alpha\) is injective and its cokernel is of rank 1 over \(\mathrm{L}(\mathrm{V},\mathrm{p}.129,\mathrm{Cor.})\) ; so we have \(d(\mathrm{L} / \mathrm{K}) = 1\) and also tr. \(\deg_{\mathbf{K}}\mathrm{L} = 1\) .
\item
  \(x\) is separable algebraic over \(\mathbf{K}\) : then \(\alpha\) is bijective (V, p.~130, Cor. 2), whence \(d(L / K) = 0\) ; clearly also tr. \(\deg_{\mathbf{K}}\mathbf{L} = 0\) .
\item
  The field L is of characteristic \(p \neq 0\) , \(x \notin K\) and \(x^{p} = a\) is in K: the cokernel C of \(\alpha\) is isomorphic to \(\Omega_{\mathrm{K}}(\mathrm{L})\) , and since \(\{x\}\) is a p-basis of L over K, the space C is of dimension 1 over L (V, p.~102, Th. 1). Since a is a p-th power in L, we have \(d_{L/P}a = 0\) and the kernel of \(\alpha\) contains the subspace R of \(U = \Omega_{\mathrm{P}}(\mathrm{K}) \otimes_{\mathrm{K}} L\) generated by \(d_{K/P}a \otimes 1\) . For each \(y \in K\) let \(\Delta(y)\) be the residue class of \(d_{K/P}y \otimes 1 \mod R\) ; then \(\Delta\) is a P-derivation of K into U/R such that \(\Delta(a) = 0\) . Prop. 5 (V, p.~101) shows that \(\Delta\) extends to a P-derivation D of L into U/R. Hence there exists an L-linear mapping \(\beta : \Omega_{\mathrm{P}}(\mathrm{L}) \to \mathrm{U}/\mathrm{R}\) such that \(D = \beta \circ d_{L/P}\) and \(\beta \circ \alpha\) is the canonical mapping of U onto U/R. This proves R to be the kernel of \(\alpha\) . Since P is perfect, we have \(P(K^{p}) = K^{p}\) , whence \(a \notin P(K^{p})\) and finally \(d_{K/P}a \neq 0\) (V, p.~103, Prop. 6). The kernel and cokernel of \(\alpha\) are thus of dimension 1, whence \(d(\mathrm{L}/\mathrm{K}) = 0\) , and we also have tr. deg \(_{K}L = 0\) .
\end{enumerate}

COROLLARY 1. --- Let L be a finitely generated extension of a field K of transcendence degree s. The vector space \(\Omega_{\mathrm{K}}(\mathrm{L})\) is of dimension \(\geqslant s\) over L, with equality if and only if L is separable over K.

Let P be the prime subfield of K. Let N be the kernel of \(\alpha\) , then by exactness of the sequence \((\mathrm{E}_{\mathrm{P},\mathrm{K},\mathrm{L}})(\mathrm{V},\mathrm{p}.128)\) and Th. 4, we have \([\Omega_{\mathrm{K}}(\mathrm{L}):\mathrm{L}] = s + [\mathrm{N}:\mathrm{L}]\) ; by V, p.~132, Cor., the extension L of K is separable if and only if N = 0. The corollary now follows.

COROLLARY 2. --- Let L be a finitely generated extension of a field K. For L to be algebraic and separable over K it is necessary and sufficient that \(\Omega_{\mathbf{K}}(\mathbf{L}) = 0\) .

This follows immediately from Cor. 1.

COROLLARY 3. --- Let K be a field of characteristic \(p \neq 0\) and L a finitely generated extension of K of transcendence degree s. If \([K:K^{p}]\) is finite, the same is true of \([L:L^{p}]\) and we have \([L:L^{p}] = p^{s} \cdot [K:K^{p}]\) .

Let P be the prime subfield of K. If \([K:K^{p}]\) is finite, then the vector space \(\Omega_{\mathrm{P}}(\mathrm{K})=\Omega_{\mathrm{K}^{\mathrm{p}}}\left(\mathrm{K}\right)\) is of finite dimension m over K, and we have \([K:K^{p}]=p^{m}(V,p.103,Th.2)\) ; the vector space \(\Omega_{\mathrm{P}}(\mathrm{K})\otimes_{\mathrm{K}}L\) is then of finite dimension m over L. Further, since K is of finite degree over \(K^{p}\) , the field \(\mathrm{K}(L^{p})\) is of finite degree over \(\mathrm{K}^{p}(L^{p})=L^{p}\) ; since the field L is a finitely generated extension of K and is algebraic over \(\mathrm{K}(L^{p})\) , it is an extension of finite degree of \(\mathrm{K}(L^{p})(\mathrm{V},p.18,Th.2)\) ; we thus conclude that L is of finite degree over \(L^{p}(V,p.10,Th.1)\) . Then \(\Omega_{\mathrm{P}}(\mathrm{L})\) is a vector space of finite dimension n over L and we have \([L:L^{p}]=p^{n}(V,p.103,Th.2)\) . By Lemma 1 (V, p.~132) the L-linear mapping \(\alpha:\Omega_{\mathrm{P}}(\mathrm{K})\otimes_{\mathrm{K}}L\to\Omega_{\mathrm{P}}(\mathrm{L})\) therefore has index n-m, whence n-m=s by Th. 4 (V, p.~133) and \(p^{n}=p^{s}\cdot p^{m}\) .

Remarks. --- 1) Let K be a field of characteristic \(p \neq 0\) and L an extension of K. We have \(\Omega_{\mathbf{K}}(\mathbf{L}) = 0\) if and only \(\mathbf{L} = \mathbf{K}(\mathbf{L}^p)\) (V, p.~103, Prop. 6). Therefore if L is finitely generated over K, then it is an algebraic and separable extension if and only if we have \(\mathbf{L} = \mathbf{K}(\mathbf{L}^p)\) . When L is not finitely generated over K, this result no longer holds generally as is shown by the case where L is the perfect closure of K.

\begin{enumerate}
\def\labelenumi{\arabic{enumi})}
\setcounter{enumi}{1}
\tightlist
\item
  Let K be a field, \(F_{1}, \ldots, F_{m}\) polynomials in \(K[X_{1}, \ldots, X_{n}]\) and L an extension of K generated by the elements \(x_{1}, \ldots, x_{n}\) . Suppose that the polynomials \(F_{1}, \ldots, F_{m}\) generate the ideal of \(K[X_{1}, \ldots, X_{n}]\) consisting of all polynomials F such that \(F(x_{1}, \ldots, x_{n}) = 0\) . From Prop. 1 (V, p.~127) and the universal property of the module of differentials (III, p.~569) we easily deduce the following result: the vector space \(\Omega_{K}(L)\) over L is generated by \(dx_{1}, \ldots, dx_{n}\) ; we have the relations
\end{enumerate}

\[
\sum_ {i = 1} ^ {n} \mathrm{D} _ {i} \mathrm{F} _ {j} (x _ {1}, \dots , x _ {n}) \cdot d x _ {i} = 0 \quad (\text { for } 1 \leqslant j \leqslant m);\tag{6}
\]

finally if \(u_1, \ldots, u_n\) are elements of \(L\) such that \(\sum_{i=1}^{n} u_i \cdot dx_i = 0\) , there exist elements \(v_1, \ldots, v_m\) of \(L\) such that \(u_i = \sum_{j=1}^{m} D_i F_j(x_1, \ldots, x_n) v_j\) for \(1 \leqslant i \leqslant n\) . Let us denote by \(r\) the rank of the matrix \((D_i F_j(x_1, \ldots, x_n))\) with \(n\) rows and \(m\) columns; let \(s\) be the transcendence degree of L over K. Then we have \([\Omega_{\mathrm{K}}(\mathrm{L}):\mathrm{L}]=n-r\) . Therefore the extension L of K is separable if and only if \(r+s=n\) (Cor. 1), and it is algebraic and separable if and only if r=n (Cor. 2).

